\begin{textsample}{NEG Dim 2 – vem\_ed\_30 – Score 3.00 – vem\_ed\_30/t157.txt}  \label{ex:f2_neg_040}
Claudia Delfino em conferência no Reino Unido Maria Claudia Delfino, professora da Fatec Praia Grande e responsável por Projetos Colaborativos Internacionais (PCIs) em inglês na área de Apoio à Internacionalização do Ensino Superior da Divisão de Extensão e Pesquisa (Depes) da Coordenadoria Geral de Ensino Superior de Graduação (CGESG) do CPS, apresentou seis trabalhos na 13th International Corpus Linguistics Conference, que ocorreu de 30 de junho a 3 de julho em Birmingham, Reino Unido.
A conferência reuniu pesquisadores de todo o mundo para discutir as mais recentes tendências e inovações na área de Linguística de Corpus.
O estudo intitulado “A large-scale key feature analysis of academic writing in the humanities” (Uma análise de características-chave em larga escala da redação acadêmica em ciências humanas), foi elaborado em coautoria com Marilisa Shimazumi, professora da Fatec Zona Leste, que também participou presencialmente no evento, e colegas de outras Instituições de Ensino Superior.
A professora Claudia apresentou ainda o pôster “Politicizing public health: The discourses around public health organizations” (Politizando a saúde pública: Os discursos em torno das organizações de saúde pública), elaborado com as colegas da Fatec Praia Grande, Tatiana Schmitz de Almeida Lopes e Fernanda Peixoto Coelho.
Veja a seguir a lista dos trabalhos apresentados pela professora.
O programa completo da conferência está disponível no site: https://www.cl2025.co.uk/programme continuação Diversidade de temas nas apresentações Os seis trabalhos desenvolvidos pela professora Claudia com colegas de Instituições de Ensino Superior brasileiras e internacionais abordaram de Inteligência Artificial a saúde pública, de vocabulário acadêmico a questões de gênero. “Shaping LLM ideology: A multi-dimensional approach” (“Moldando a ideologia de Large Language Models: uma abordagem multidimensional”) Autores: Tony Berber Sardinha (PUC-SP), Anderson Avila (Institut Nacional de la Recherche Scientifique / Université du Québec en Outaouais), Maria Claudia Delfino (Fatec Praia Grande), Hazem Amamou (Institut Nacional de la Recherche Scientifique / Université du Québec en Outaouais), Rogerio Yamada (PUC-SP) e Ru-bing Chen (The Hong Kong Politechnic) “A large-scale key feature analysis of academic writing in the humanities” (“Uma análise de características-chave em larga escala da redação acadêmica em ciências humanas”) Autores: Tony Berber Sardinha (PUC-SP), Deise Dutra (UFMG), Maria Claudia Delfino (Fatec Praia Grande), Ana Bocorny (UFRGS), Marilisa Shimazumi (Fatec Zona Leste), Carlos Kauffmann (PUC-SP), Luciana Dias Macedo (UFMG) e Ana Clara Taborda (UFMG) continuação “Artificial intelligence in songwriting: A lexical multi-dimensional approach” (Inteligência artificial na composição de músicas: Uma abordagem lexical multidimensional) Autores: Maria Claudia Delfino e Tony Berber Sardinha (PUC-SP) “Online discourses of ‘Gender Critical’ feminists: A keyword co-occurrence analysis” (“Discursos on-line de feministas ´ críticas de gênero ´: Uma análise de co-ocorrência de palavras-chave”) Autores: Mark McGlashan (University of Liverpool), Isobelle Clarke (Lancaster University), Tony Berber-Sardinha, Maria Claudia Delfino e Mirella Whiteman (PUC-SP) “The role of academic vocabulary in shaping Sustainable Development Goal (SDG) research papers: A lexical multidimensional perspective” (“O papel do vocabulário acadêmico na formação de artigos de pesquisa sobre os Objetivos de Desenvolvimento Sustentável (ODS): Uma perspectiva lexical multidimensional”) Autores: Paula Pinto (Unesp), Tony Berber Sardinha (PUC-SP), Denis Owa (PUC-SP), Maria Claudia Delfino (Fatec Praia Grande) e Simone Resende (Faculdade Cultura Inglesa) “Politicizing public health: The discourses around public health organizations” (“Politizando a saúde pública: Os discursos em torno das organizações de saúde pública”) Autores: Tatiana Schmitz de Almeida Lopes, Maria Claudia Delfino e Fernanda Peixoto Coelho (Fatec Praia Grande)

% matched lemmas: 
\end{textsample}
